% formal/flow-graph-real-algorithm.tex
%
% Status: agent-written proof-development surface, unapproved by Jorn.
% Purpose: define the idealized real-number flow-graph/tube algorithm and
% isolate the correctness theorem needed before thesis prose is written.
%
% Primary recovery sources:
% - crates/symplectic/src/algorithms/flow_graph/tube-algorithm-raw-jorn-2026-05-04.md
% - crates/symplectic/src/algorithms/flow_graph/tube-algorithm-legacy-source-note.md
% - git show 25dd8b9acb8aeaaa6aa3abd80fc6d95db00c4747:formal/tube-algorithm.tex
% - formal/reeb-orbit-recovery.tex
% - formal/hk2017-qp-core.tex

\begin{unverified}
\begin{fact}[Simple capacity minimizer]\label{fact:fg-simple-capacity-minimizer}
% Source: thesis theorem thm:generalized-reeb-simple-minimizer, which records
% the HK2017 simple-minimizer argument in thesis conventions.
For every full-dimensional convex polytope
\(K\subset\mathbb R^4\) with \(0\in\operatorname{int}K\),
\(c_{\mathrm{EHZ}}(K)\) is achieved by a simple closed generalized Reeb
orbit.
\end{fact}
\end{unverified}

\begin{unverified}
\begin{definition}[Flow-graph nondegenerate facet presentation]%
\label{def:fg-nondegenerate-facet-presentation}
Let
\[
  K=\{x\in\mathbb R^4:\langle a_i,x\rangle\leq 1,
      \ i=1,\ldots,F\}
\]
be a full-dimensional convex polytope with \(0\in\operatorname{int}K\).
Write \(F_i=K\cap\{\langle a_i,x\rangle=1\}\) and
\(R_i=2J_0a_i\).
Also write \(H_i=\{x:\langle a_i,x\rangle=1\}\).
Assume each \(F_i\) is a genuine facet of \(K\).

A \emph{physical facet transition} \(F_i\to F_j\), for distinct facet indices
\(i\ne j\), exists if there are \(x\in F_i\cap F_j\) and
\(\varepsilon>0\) such that
\[
  x-tR_i\in K\cap H_i
  \quad\text{and}\quad
  x+tR_j\in K\cap H_j
  \qquad(0\le t\le\varepsilon).
\]
This is the local breakpoint condition for a trajectory segment on \(F_i\)
before the breakpoint and on \(F_j\) after the breakpoint.  The definition
intentionally allows the breakpoint to lie on additional facets; it is not
restricted to relative-interior points of two-faces.

The facet presentation is \emph{flow-graph nondegenerate} if:
\begin{enumerate}
\item whenever \(i\ne j\) and \(F_i\cap F_j\ne\emptyset\),
  \(\omega_0(a_i,a_j)\ne0\);
\item every subset of at most four vectors among
  \(a_1,\ldots,a_F\) is linearly independent.
\end{enumerate}
\end{definition}
\end{unverified}

\begin{unverified}
\begin{definition}[Flow-graph transition sign]%
\label{def:fg-transition-sign}
For an ordered triple of facet indices \((p,i,n)\), say that
\((p,i,n)\) has \emph{positive flow-graph transition sign} if
\[
  \langle R_p,a_i\rangle>0
  \qquad\text{and}\qquad
  \langle R_i,a_n\rangle>0.
\]
\end{definition}
\end{unverified}

\begin{unverified}
\begin{lemma}[Local transition regularity gives positive transition signs]%
\label{lem:fg-local-transition-regularity-positive-sign}
Assume the facet presentation is flow-graph nondegenerate.  If a physical
facet transition \(F_i\to F_j\) exists, then
\[
  \langle R_i,a_j\rangle>0 .
\]
Consequently, if a strict Reeb trajectory has positive dwell time on
successive facets \(F_p,F_i,F_n\), then the triple \((p,i,n)\) has positive
flow-graph transition sign.
\end{lemma}

\begin{proof}
Let \(x\in F_i\cap F_j\) and \(\varepsilon>0\) witness the physical
transition.  Since \(x-tR_i\in K\) for \(0\le t\le\varepsilon\), the
\(j\)-th defining inequality gives
\[
  1-t\langle a_j,R_i\rangle
  =
  \langle a_j,x-tR_i\rangle
  \le 1 .
\]
Thus \(\langle a_j,R_i\rangle\ge0\).  The local nonzero-\(\omega_0\)
condition in Definition~\ref{def:fg-nondegenerate-facet-presentation} excludes
equality, because a physical transition has \(x\in F_i\cap F_j\), so the
nonempty facet-pair condition applies, and
\[
  \langle a_j,R_i\rangle=2\omega_0(a_i,a_j).
\]
Hence \(\langle R_i,a_j\rangle=\langle a_j,R_i\rangle>0\).

For a strict trajectory with positive dwell time on
\(F_p,F_i,F_n\), the breakpoint from \(F_p\) to \(F_i\) is a physical
transition \(F_p\to F_i\), and the breakpoint from \(F_i\) to \(F_n\) is a
physical transition \(F_i\to F_n\).  Applying the first part to these two
transitions gives
\(\langle R_p,a_i\rangle>0\) and \(\langle R_i,a_n\rangle>0\), which is
Definition~\ref{def:fg-transition-sign}.
\end{proof}
\end{unverified}

\begin{unverified}
\begin{definition}[Simple cyclic facet word]\label{def:fg-simple-cyclic-word}
A \emph{simple cyclic facet word} is a finite cyclic equivalence class
\[
  [s_1,\ldots,s_k],
  \qquad k\ge2,
\]
where cyclic rotations are identified and the listed facet indices are
distinct.  In formulas we write a representative \((s_1,\ldots,s_k)\), but
the word is the cyclic class.
\end{definition}
\end{unverified}

\begin{unverified}
\begin{definition}[Closed tube search data]\label{def:fg-closed-tube-search-data}
Let \(s=(s_1,\ldots,s_{k+2})\) be a facet sequence.  A strict tube for this
sequence is a Reeb trajectory segment whose breakpoints lie successively in
\[
  F_{s_1}\cap F_{s_2},\;
  F_{s_2}\cap F_{s_3},\ldots,
  F_{s_{k+1}}\cap F_{s_{k+2}},
\]
and whose \(r\)-th segment flows along \(R_{s_{r+1}}\) for time
\(\tau_r>0\).

For each facet pair \((i,j)\) used below, choose an affine coordinate chart
\[
  \psi_{ij}\colon\mathbb R^2\to H_i\cap H_j .
\]
The corresponding local face polygon is
\[
  P_{ij}=\{u\in\mathbb R^2:\psi_{ij}(u)\in K\}.
\]
The algorithm works with the closed search data associated to the inner
segments of this strict object.  It consists of:
\begin{itemize}
\item a closed start polygon \(P_{\mathrm{start}}\) in an affine coordinate
  chart for \(H_{s_1}\cap H_{s_2}\), contained in \(P_{s_1s_2}\);
\item a closed end polygon \(P_{\mathrm{end}}\) in an affine coordinate chart
  for \(H_{s_{k+1}}\cap H_{s_{k+2}}\), contained in
  \(P_{s_{k+1}s_{k+2}}\);
\item an affine map \(\Phi\colon\mathbb R^2\to\mathbb R^2\) with
  \(\Phi(P_{\mathrm{start}})\subseteq P_{\mathrm{end}}\);
\item an affine action function \(A\colon\mathbb R^2\to\mathbb R\).
\end{itemize}
It represents the closed inner candidate data if, for every
\(u_1\in P_{\mathrm{start}}\), the breakpoints
\[
  x_r\in F_{s_r}\cap F_{s_{r+1}},
  \qquad r=1,\ldots,k+1,
\]
are determined recursively by
\[
  x_1=\psi_{s_1s_2}(u_1),
  \qquad
  x_{r+1}=x_r+\tau_rR_{s_{r+1}},
  \qquad
  \tau_r\ge0,
\]
and satisfy
\[
  \psi_{s_{k+1}s_{k+2}}^{-1}(x_{k+1})=\Phi(u_1),
  \qquad
  A(u_1)=\sum_{r=1}^{k}\tau_r .
\]
Here \(\psi_{ij}\) denotes the chosen affine chart for
\(H_i\cap H_j\).  The point \(u_1\) represents a strict inner tube exactly
when all reconstructed segment times are positive.  For an open tube piece,
raw one-sided membership also depends on compatible neighboring pieces; for a
closed word these neighboring pieces are supplied by the same cyclic word.
\end{definition}
\end{unverified}

\begin{unverified}
\begin{definition}[Primitive flow-graph tube]%
\label{def:fg-primitive-tube}
Assume \((p,i,n)\) has positive flow-graph transition sign.  For
\(x\in F_p\cap F_i\), define
\[
  \tau_{pin}(x)
  =
  \frac{1-\langle a_n,x\rangle}{\langle a_n,R_i\rangle},
  \qquad
  \varphi_{pin}(x)=x+\tau_{pin}(x)R_i .
\]
The primitive closed search data for \((p,i,n)\) is obtained by choosing
facet-pair charts for \(H_p\cap H_i\) and \(H_i\cap H_n\), restricting
the start coordinate to the closed set where \(x\in K\),
\(\varphi_{pin}(x)\in K\), and \(\tau_{pin}(x)\geq0\), and by using the
coordinate expression of \(\varphi_{pin}\) as the affine map and
\(\tau_{pin}\) as the action function.
\end{definition}
\end{unverified}

\begin{unverified}
\begin{lemma}[Primitive tubes are affine closed search data]%
\label{lem:fg-primitive-tubes-affine}
For every triple \((p,i,n)\) with positive flow-graph transition sign, the
primitive construction of Definition~\ref{def:fg-primitive-tube} defines
closed tube search data in the sense of
Definition~\ref{def:fg-closed-tube-search-data}.  Its strict points are
exactly the points with \(\tau_{pin}>0\) for the inner primitive segment.
The affine map is an affine bijection between the two ambient facet-pair
planes.
\end{lemma}

\begin{proof}
The denominator \(\langle a_n,R_i\rangle\) is positive by
Definition~\ref{def:fg-transition-sign}, so \(\tau_{pin}\) is an affine
function on \(F_p\cap F_i\).  The map
\(\varphi_{pin}(x)=x+\tau_{pin}(x)R_i\) is affine and lands in the hyperplane
\(\langle a_n,\cdot\rangle=1\).  Since \(R_i\) is tangent to \(F_i\), it also
lands in the hyperplane \(\langle a_i,\cdot\rangle=1\).  Pulling back the
halfspace inequalities of \(K\), the endpoint face inequalities, and
\(\tau_{pin}\ge0\) gives a closed polygonal start domain.

The inverse affine map is obtained by flowing backwards along \(R_i\).  For
\(y\in H_i\cap H_n\), put
\[
  \sigma_{inp}(y)=
  \frac{1-\langle a_p,y\rangle}{\langle a_p,R_i\rangle}.
\]
The denominator is nonzero because
\[
  \langle a_p,R_i\rangle
  =-\langle R_p,a_i\rangle,
\]
and the last quantity is nonzero under the positive transition sign
hypothesis.  The point \(y+\sigma_{inp}(y)R_i\) lies in \(H_i\cap H_p\), and
if \(y=x+\tau_{pin}(x)R_i\) with \(x\in H_p\cap H_i\), then
\(\sigma_{inp}(y)=-\tau_{pin}(x)\).  The same computation in the other
direction gives the inverse relation on the two affine planes.

For any start-domain point, the endpoint
\(\varphi_{pin}(x)\) lies in \(K\) by definition of the start domain.  The
whole segment lies in \(K\) because every defining inequality of \(K\) is
affine along the segment and both endpoints are in \(K\).  The whole segment
lies in \(H_i\) because \(R_i\) is tangent to \(H_i\).  Thus the primitive
affine data represent exactly the closed primitive inner candidates, with
action \(\tau_{pin}\).  The strict inner segment condition is exactly
\(\tau_{pin}>0\).
\end{proof}
\end{unverified}

\begin{unverified}
\begin{lemma}[Gluing closed tube search data]\label{lem:fg-tube-gluing}
If closed tube search data for \(s=(s_1,\ldots,s_{k+2})\) and
\(t=(s_{k+1},s_{k+2},\ldots,s_{k+\ell+2})\) have matching middle
facet-pair plane \(H_{s_{k+1}}\cap H_{s_{k+2}}\), write their data as
\[
  (P_1,\widehat P_1,\Phi_1,A_1),
  \qquad
  (P_2,\widehat P_2,\Phi_2,A_2).
\]
The concatenated start polygon is
\[
  P_1\cap\Phi_1^{-1}(P_2),
\]
the concatenated end polygon is
\(\Phi_2(P_2\cap\Phi_1(P_1))\), the concatenated affine map is
\(\Phi_2\circ\Phi_1\), and these data give the closed tube search data for the
concatenated facet sequence.

The action function of the concatenation is the sum of the first action
function and the pullback of the second action function.
\end{lemma}

\begin{proof}
Both pieces use the same affine coordinates on the shared facet-pair plane.
A point in \(P_1\) survives exactly when \(\Phi_1(u)\in P_2\), which is the
preimage condition above.  For such a point, the reconstructed endpoint of the
first piece is the reconstructed start point of the second piece.  Hence
concatenating the two breakpoint lists and segment-time lists gives a valid
candidate for the concatenated sequence.  Its endpoint coordinate is
\(\Phi_2(\Phi_1(u))\), and its action is
\(A_1(u)+A_2(\Phi_1(u))\).

Conversely, any candidate for the concatenated sequence restricts to a
candidate for the first sequence and to a candidate for the second sequence,
with the shared breakpoint giving the same middle coordinate.  Therefore its
start coordinate lies in \(P_1\cap\Phi_1^{-1}(P_2)\), and its endpoint and
action are given by the displayed formulas.  Since affine maps compose and
affine functions are closed under pullback and addition, the glued data are
again closed tube search data.
\end{proof}
\end{unverified}

\begin{unverified}
\begin{lemma}[Empty subtubes stay empty under gluing]%
\label{lem:fg-empty-subtube-pruning}
In the notation of Lemma~\ref{lem:fg-tube-gluing}, if
\[
  P_1\cap\Phi_1^{-1}(P_2)=\emptyset,
\]
then the concatenated closed tube search data has empty start polygon and
there is no closed search-domain candidate for the concatenated facet
sequence.  More generally, any iterated gluing of closed tube search data has
empty start polygon once one of its gluing stages has empty start polygon.
\end{lemma}

\begin{proof}
The first statement is the start-polygon formula from
Lemma~\ref{lem:fg-tube-gluing}.  A candidate for the concatenated sequence
would have to restrict to a candidate for both pieces with matching middle
coordinate, which is exactly the condition
\(u\in P_1\cap\Phi_1^{-1}(P_2)\).  If this set is empty, no such candidate
exists.

For iterated gluing, apply the first statement at the first empty stage.  Every
later glued tube has a start polygon obtained by imposing additional preimage
conditions on an already empty start polygon, so it remains empty.
\end{proof}
\end{unverified}

\begin{unverified}
\begin{lemma}[Closed tube fixed points]\label{lem:fg-closed-tube-fixed-points}
Let \(w=(s_1,\ldots,s_k)\) be a cyclic facet word and let the closed tube
search data for
\[
  (s_1,\ldots,s_k,s_1,s_2)
\]
have start polygon \(P\), affine map \(\Phi(u)=Mu+b\), and action function
\(A\).  The closed search-domain candidates for this word are exactly the
points \(u\in P\) satisfying \(\Phi(u)=u\).  Such a fixed point represents a
strict closed Reeb trajectory with cyclic word \(w\) exactly when every
reconstructed segment time is positive.

If \(I-M\) is invertible, there is at most one fixed point.  If \(I-M\) is
singular, the fixed-point set can be empty, nonisolated, or positive
dimensional; algebraically it is the intersection of \(P\) with the affine
solution set of \((I-M)u=b\).
\end{lemma}

\begin{proof}
A fixed point means that the concatenated tube starts and ends at the same
point of \(F_{s_1}\cap F_{s_2}\).  This is precisely the closed search-domain
condition.  The distinction between a closed search-domain point and the strict
orbit represented by the displayed word is the strict positivity of all segment
times, by Definition~\ref{def:fg-closed-tube-search-data}.  The final
statements are the elementary linear algebra of \((I-M)u=b\).
\end{proof}
\end{unverified}

\begin{unverified}
\begin{lemma}[Flow-graph action normalization]\label{lem:fg-action-normalization}
For a strict tube segment on \(F_i\) flowing for time \(\tau\) along
\(R_i=2J_0a_i\), the symplectic action contribution is \(\tau\).  Hence the
action function of a glued tube is the sum of the reconstructed segment times.
\end{lemma}

\begin{proof}
This is the contact-normalized convention used for the generalized Reeb orbit
model in this formal tree.  With the standard primitive
\(\lambda_0=\frac12\omega_0(x,\cdot)\), the point of the segment lies on
\(\langle a_i,x\rangle=1\), and the convention \(R_i=2J_0a_i\) gives
\(\lambda_0(R_i)=1\).  Integrating \(\lambda_0(\dot\gamma)\) over a segment of
duration \(\tau\) gives \(\tau\).  Additivity over glued segments gives the
second statement.
\end{proof}
\end{unverified}

\begin{unverified}
\begin{lemma}[Action cutoffs preserve below-cutoff fixed points]%
\label{lem:fg-action-cutoff-preserves-below-cutoff-fixed-points}
Let a closed tube search datum for a cyclic word have start polygon \(P\),
affine map \(\Phi\), and action function \(A\).  For \(B\in\mathbb R\), let
\[
  P_B=P\cap\{u:A(u)\le B\}.
\]
Then the fixed points of \(\Phi\) in \(P_B\) are exactly the fixed points of
\(\Phi\) in \(P\) whose action is at most \(B\).
\end{lemma}

\begin{proof}
This is immediate from the definition of \(P_B\).  A point \(u\) is in \(P_B\)
and satisfies \(\Phi(u)=u\) if and only if \(u\in P\), \(A(u)\le B\), and
\(\Phi(u)=u\).
\end{proof}
\end{unverified}

\begin{unverified}
\begin{lemma}[Dynamic action cutoffs preserve retained output]%
\label{lem:fg-dynamic-action-cutoffs-preserve-retained-output}
Fix a finite collection of closed tube search data and a threshold
\(\delta\ge0\).  During an exhaustive search, suppose that after at least one
strict closed trajectory has been found, the search restricts every later tube
by the cutoff \(B=c_{\mathrm{known}}+\delta\), where \(c_{\mathrm{known}}\) is
the least action among the strict trajectories already found.  Then these
dynamic action cutoffs do not change the final minimum action and do not remove
any strict trajectory whose action is at most the final minimum plus
\(\delta\).
\end{lemma}

\begin{proof}
Let \(c\) be the minimum action among all strict trajectories in the uncut
finite search.  At every cutoff step, \(c\le c_{\mathrm{known}}\), because
\(c_{\mathrm{known}}\) is the action of an already found strict trajectory.
Therefore
\[
  c+\delta\le c_{\mathrm{known}}+\delta=B .
\]
By Lemma~\ref{lem:fg-action-cutoff-preserves-below-cutoff-fixed-points}, the
cutoff keeps all fixed points with action at most \(B\).  In particular it
keeps every trajectory with action at most \(c+\delta\), including every
minimum-action trajectory that appears later in the enumeration.  Hence the
final minimum and the retained set below final minimum plus \(\delta\) agree
with the uncut exhaustive search.
\end{proof}
\end{unverified}

\begin{unverified}
\begin{lemma}[Nonpositive fixed sets contain no strict orbit]%
\label{lem:fg-nonpositive-fixed-set-no-strict-orbit}
Let \(w=[s_1,\ldots,s_k]\) be a cyclic facet word and let the closed tube
search data for
\[
  (s_1,\ldots,s_k,s_1,s_2)
\]
have start polygon \(P\), affine map \(\Phi\), and action function \(A\).  Let
\[
  C=\{u\in P:\Phi(u)=u\}
\]
be the closed fixed-point set.  If \(A(u)\le0\) for every \(u\in C\), then
this closed tube has no strict represented orbit with cyclic word \(w\).

In particular, if \(C\) is a compact convex polygon and \(A(v)\le0\) for every
vertex \(v\) of \(C\), then the same conclusion holds.
\end{lemma}

\begin{proof}
For a closed search-domain fixed point, the action \(A(u)\) is the sum of the
reconstructed segment times by
Definition~\ref{def:fg-closed-tube-search-data}.  Each reconstructed segment
time is nonnegative.  If \(u\) represented a strict orbit for the displayed
word, then every segment time would be positive, and hence \(A(u)>0\).  This
contradicts \(A(u)\le0\).

For the final statement, an affine function on a compact convex polygon attains
its maximum at a vertex.  If all vertex values are nonpositive, then
\(A(u)\le0\) for every \(u\in C\), so the first part applies.
\end{proof}
\end{unverified}

\begin{unverified}
\begin{lemma}[Simple support orbits give flow-graph fixed points]%
\label{lem:fg-simple-support-orbit-represented}
Let \(\gamma\) be a simple closed Reeb orbit whose support word
\([s_1,\ldots,s_k]\) has positive dwell time on every listed facet.  If every
adjacent triple of this support word has positive flow-graph transition sign,
then \(\gamma\) determines a strict fixed point of the closed tube search data
for \((s_1,\ldots,s_k,s_1,s_2)\), and the flow-graph action of that fixed point
equals the Reeb action of \(\gamma\).
\end{lemma}

\begin{proof}
Choose the breakpoint at which \(\gamma\) enters
\(F_{s_1}\cap F_{s_2}\) as the start point.  Since the orbit has positive dwell
time on every listed facet, its successive breakpoints lie on the two-faces
\(F_{s_r}\cap F_{s_{r+1}}\), and each segment is obtained by flowing along
\(R_{s_{r+1}}\) for a positive time.  The primitive construction recovers the
same endpoint and segment time from each start point.  Gluing the primitives
therefore reconstructs the whole closed orbit as a closed tube search-domain
fixed point, and the positivity of the dwell times makes it strict.  The action
identity follows from Lemma~\ref{lem:fg-action-normalization}.
\end{proof}
\end{unverified}

\begin{unverified}
\begin{algorithm}[Idealized exhaustive flow-graph tube search]%
\label{alg:fg-real-exhaustive-search}
Input: a polytope \(K\) with a flow-graph nondegenerate facet presentation in
the sense of Definition~\ref{def:fg-nondegenerate-facet-presentation}, and a
threshold \(\delta\ge0\).

\begin{enumerate}
\item Enumerate every simple cyclic facet word \(w=[s_1,\ldots,s_k]\) in the
  sense of Definition~\ref{def:fg-simple-cyclic-word}.
\item For each word, form the closed sequence
  \((s_1,\ldots,s_k,s_1,s_2)\).  If one of its adjacent facet pairs has empty
  intersection, or if one of its triples does not have positive flow-graph
  transition sign, continue with no orbit for this word.
\item Build the closed tube search data by gluing primitive tubes.  If the
  resulting start polygon is empty, continue with no orbit for this word.
\item Solve \(\Phi_w(u)=u\).  If the fixed-point equation is singular and
  \(|w|\le4\), continue with no orbit for this word by
  Lemma~\ref{lem:fg-short-words-no-positive-orbit}.  If it is singular and
  \(|w|\ge5\), stop and report that the input is outside this algorithm's
  regular case.
\item Record a fixed point only when it lies in the closed start polygon and
  every reconstructed segment time is positive.
\item If no fixed point was recorded, return that no represented strict closed
  trajectory was found.  Otherwise return the minimum recorded action and all
  recorded strict closed trajectories with action at most this minimum plus
  \(\delta\).
\end{enumerate}
\end{algorithm}
\end{unverified}

\begin{unverified}
\begin{definition}[Finite-orbit regularity]\label{def:fg-finite-orbit-regularity}
Assume the facet presentation is flow-graph nondegenerate in the sense of
Definition~\ref{def:fg-nondegenerate-facet-presentation}.  Build the closed
tube search data for every simple cyclic word that survives the adjacent
facet-pair, transition-sign, and nonempty-domain steps of
Algorithm~\ref{alg:fg-real-exhaustive-search}.
For such a word \(w\), write the affine tube map as \(\Phi_w(u)=M_wu+b_w\).

The run is \emph{finite-orbit regular} if, for every enumerated word whose
closed search domain is nonempty and whose length is at least five,
\[
  \det(I-M_w)\ne0.
\]
Words of length at most four are handled by the linear-dependence argument in
Lemma~\ref{lem:fg-short-words-no-positive-orbit}.
\end{definition}
\end{unverified}

\begin{unverified}
\begin{lemma}[Short simple words have no positive-time fixed point]%
\label{lem:fg-short-words-no-positive-orbit}
Assume the dual vertices satisfy the linear-independence condition in
Definition~\ref{def:fg-nondegenerate-facet-presentation}.  If a simple cyclic word
has length at most four, then it has no strict positive-time closed Reeb
trajectory.
\end{lemma}

\begin{proof}
For a strict positive-time closed trajectory with cyclic word
\((s_1,\ldots,s_k)\), closure gives
\[
  \sum_{r=1}^k \tau_r R_{s_r}=0
  \qquad\text{with all }\tau_r>0.
\]
Since \(R_i=2J_0a_i\) and \(J_0\) is invertible, this is equivalent to a
positive linear dependence among \(a_{s_1},\ldots,a_{s_k}\).  If \(k\le4\),
this contradicts the assumed linear independence of every subset of at most
four dual vertices.
\end{proof}
\end{unverified}

\begin{unverified}
\begin{lemma}[Primitive linear maps are symplectic]%
\label{lem:fg-primitive-linear-maps-symplectic}
Let \((p,i,n)\) be a primitive triple with
\(\langle a_n,R_i\rangle\ne0\) and \(\omega_0(a_p,a_i)\ne0\).  Put
\[
  V_{pi}=\ker a_p\cap\ker a_i,
  \qquad
  V_{in}=\ker a_i\cap\ker a_n .
\]
The linear part
\[
  D_{pin}(\xi)=
  \xi
  -
  \frac{\langle a_n,\xi\rangle}{\langle a_n,R_i\rangle}R_i
\]
maps \(V_{pi}\) to \(V_{in}\) and preserves the restriction of \(\omega_0\):
\[
  \omega_0(D_{pin}\xi,D_{pin}\eta)=\omega_0(\xi,\eta)
  \qquad(\xi,\eta\in V_{pi}).
\]
Moreover \(V_{pi}\) is a two-dimensional symplectic vector space.  Hence every
closed return map built from such primitive steps has determinant \(1\) on its
two-dimensional start plane.
\end{lemma}

\begin{proof}
For \(\xi\in V_{pi}\), the formula gives
\[
  \langle a_i,D_{pin}\xi\rangle=\langle a_i,\xi\rangle=0
\]
because \(R_i=2J_0a_i\) is tangent to \(H_i\), and
\[
  \langle a_n,D_{pin}\xi\rangle
  =
  \langle a_n,\xi\rangle
  -
  \frac{\langle a_n,\xi\rangle}{\langle a_n,R_i\rangle}
  \langle a_n,R_i\rangle
  =0.
\]
Thus \(D_{pin}\xi\in V_{in}\).

For \(\xi,\eta\in V_{pi}\), both vectors lie in \(\ker a_i\).  Therefore the
cross terms with \(R_i\) vanish:
\[
  \omega_0(R_i,\eta)=0,\qquad \omega_0(\xi,R_i)=0 .
\]
Also \(\omega_0(R_i,R_i)=0\).  Expanding the displayed formula for
\(D_{pin}\) gives
\[
  \omega_0(D_{pin}\xi,D_{pin}\eta)=\omega_0(\xi,\eta).
\]

The tangent space \(V_{pi}\) has dimension two because \(a_p\) and \(a_i\) are
linearly independent.  A vector lies in \(V_{pi}\) exactly when it is
\(\omega_0\)-orthogonal to both \(R_p\) and \(R_i\).  Hence the symplectic
orthogonal of \(V_{pi}\) is \(\operatorname{span}(R_p,R_i)\).  The
restriction of \(\omega_0\) to \(V_{pi}\) is degenerate exactly when
\(V_{pi}\cap\operatorname{span}(R_p,R_i)\) contains a nonzero vector.

Write such a vector as \(\alpha R_p+\beta R_i\).  The conditions
\(\alpha R_p+\beta R_i\in V_{pi}\) are
\[
  \beta\langle a_p,R_i\rangle=0,
  \qquad
  \alpha\langle a_i,R_p\rangle=0 .
\]
The hypothesis \(\omega_0(a_p,a_i)\ne0\) makes these two pairings nonzero, so
\(\alpha=\beta=0\).  Thus \(V_{pi}\) is symplectic.  A closed return map is a
composition of primitive linear maps and starts and ends on the same
symplectic two-plane.  It preserves the symplectic area form on that plane, so
its determinant is \(1\).
\end{proof}
\end{unverified}

\begin{unverified}
\begin{lemma}[Repeated section contraction]%
\label{lem:fg-return-determinant-contraction}
Let \(x,A,B,y\) be facet labels for which the primitive denominators below are
nonzero.  In the linear primitive-transition notation of
Lemma~\ref{lem:fg-primitive-linear-maps-symplectic},
\[
  D_{BAy}\,D_{ABA}\,D_{xAB} = D_{xAy}.
\]
Thus, at the level of the linear return map, replacing a consecutive subword
\(x,A,B,A,y\) by \(x,A,y\) does not change the surrounding composed map.
\end{lemma}

\begin{proof}
For \(\xi\in\ker a_x\cap\ker a_A\), write
\[
  \eta = D_{xAB}\xi
  =
  \xi
  -
  \frac{\langle a_B,\xi\rangle}{\langle a_B,R_A\rangle}R_A .
\]
Then \(\eta\in\ker a_A\cap\ker a_B\).  Since the next primitive step is
\((A,B,A)\) and \(\eta\in\ker a_A\), its formula gives
\(D_{ABA}\eta=\eta\).

Applying the last step gives
\[
\begin{aligned}
  D_{BAy}\eta
  &=
  \eta
  -
  \frac{\langle a_y,\eta\rangle}{\langle a_y,R_A\rangle}R_A  \\
  &=
  \xi
  -
  \frac{\langle a_B,\xi\rangle}{\langle a_B,R_A\rangle}R_A
  -
  \frac{
    \langle a_y,\xi\rangle
    -
    \frac{\langle a_B,\xi\rangle}{\langle a_B,R_A\rangle}
    \langle a_y,R_A\rangle}
  {\langle a_y,R_A\rangle}R_A  \\
  &=
  \xi
  -
  \frac{\langle a_y,\xi\rangle}{\langle a_y,R_A\rangle}R_A
  =
  D_{xAy}\xi .
\end{aligned}
\]
\end{proof}
\end{unverified}

\begin{unverified}
\begin{lemma}[Long-word determinants are generically nonzero]%
\label{lem:fg-long-word-determinants-open-dense}
Fix a word length \(k\ge5\).  Let \(D_k\subset(\mathbb R^4)^k\) be the set of
tuples \(b=(b_1,\ldots,b_k)\) for which
\[
  \langle b_{r+2},2J_0b_{r+1}\rangle\ne0,
  \qquad r=1,\ldots,k,
\]
where indices are read cyclically.  On \(D_k\), the linear return map \(L_k\)
on
\[
  \ker b_1\cap\ker b_2
\]
for the cyclic primitive sequence \((b_1,\ldots,b_k,b_1,b_2)\) depends
rationally on the listed dual vertices.  Therefore, on each
coordinate-minor chart for the varying facet-pair planes, the determinant
\[
  \det(I-L_k)
\]
is a rational function of the listed dual vertices.  This rational function is
not the zero rational function.  Hence its nonzero locus is open dense in
\(D_k\).
\end{lemma}

\begin{proof}
The set \(D_k\) is open because it is defined by finitely many nonvanishing
polynomial inequalities.  The denominator condition implies that every
adjacent cyclic pair has nonzero symplectic pairing, after reindexing the
displayed inequalities.  Hence the two covectors in each adjacent pair are
linearly independent, so their common kernel is a two-dimensional plane.

For the primitive step from
\(\ker b_r\cap\ker b_{r+1}\) to
\(\ker b_{r+1}\cap\ker b_{r+2}\), the linear map is
\[
  \xi\mapsto
  \xi
  -
  \frac{\langle b_{r+2},\xi\rangle}
       {\langle b_{r+2},2J_0b_{r+1}\rangle}
  2J_0b_{r+1} .
\]
Its coefficients are rational functions of the involved dual vertices on the
open set where the denominator is nonzero.  On a coordinate patch where a
fixed \(2\times2\) minor of the two covectors defining each facet-pair plane
is nonzero, choose the two free coordinate directions as a basis of that plane
after solving the remaining two coordinates.  These basis vectors depend
rationally on the listed dual vertices.  In these bases the primitive matrices
and their finite composition have rational entries.  The determinant
\(\det(I-L_k)\) is independent of the chosen basis of the start plane, so
these local rational expressions agree on overlaps.

It remains to prove that this rational function is not identically zero.
Because different simple words of the same length differ only by relabeling
the variables, it is enough to give one nonzero specialization for each
length \(k\).  Also, changing the cyclic starting point conjugates the closed
return map by an invertible primitive transition map, so it does not change
\(\det(I-L_k)\).

For this nonvanishing check, regard the determinant as a rational expression
in formal vector variables \(b_1,\ldots,b_k\).  The expression is defined at
any vector tuple where the denominators and coordinate-minor choices used above
are nonzero.  If the rational expression were zero, every such evaluation would
be zero.  Therefore it is enough to give one algebraic vector tuple for each
length with nonzero determinant value.

Use the five vectors
\[
\begin{aligned}
 A&=(2,3,-4,-3),&
 B&=(0,1,2,-1),&
 C&=(-4,-3,0,-1),\\
 D&=(4,-1,1,1),&
 E&=(-4,2,-4,-4).
\end{aligned}
\]
The pairings \(\langle q,2J_0p\rangle\) for unordered pairs are
\[
\begin{array}{c|cccccccccc}
pq&AB&AC&AD&AE&BC&BD&BE&CD&CE&DE\\ \hline
\langle q,2J_0p\rangle
&8&-56&36&-60&8&-16&12&-16&60&-20 .
\end{array}
\]
Thus all primitive denominators used below are nonzero.

For odd \(k\ge5\), write \(k=2m+3\) with \(m\ge1\) and specialize the listed
variables cyclically as
\[
  (A,B)^m,C,D,E .
\]
If \(m>1\), rotate the cyclic word to expose a subword \(E,A,B,A,B\) and apply
Lemma~\ref{lem:fg-return-determinant-contraction} with
\((x,A,B,y)=(E,A,B,B)\).  This removes one repeated \(A,B\) pair and preserves
\(\det(I-L)\).  Repeating this contraction reduces the odd case to
\((A,B,C,D,E)\).  A direct exact computation from the displayed
primitive-transition formula gives the basis-independent value
\[
  \det(I-L_{ABCDE})=-\frac{10643}{600}\ne0 .
\]

For even \(k\ge6\), write \(k=2m+4\) with \(m\ge1\) and specialize the listed
variables cyclically as
\[
  (A,B)^m,A,C,D,E .
\]
If \(m>1\), the same \(E,A,B,A,B\) contraction reduces to the case
\((A,B,A,C,D,E)\).  In that case, rotate to expose the cyclic subword
\(E,A,B,A,C\) and apply
Lemma~\ref{lem:fg-return-determinant-contraction} with
\((x,A,B,y)=(E,A,B,C)\).  This reduces the even case to \((A,C,D,E)\).  A
direct exact computation from the displayed primitive-transition formula gives
the basis-independent value
\[
  \det(I-L_{ACDE})=\frac{1}{168}\ne0 .
\]
% 2026-06-22 Codex rechecked the displayed pairing table and the two displayed
% determinant values by an exact rational implementation of the primitive
% linear maps.  The same check gives
% \(\det(I-L_{ABACDE})=\det(I-L_{ACDE})=1/168\), matching the even-case
% contraction step.

Thus the determinant is a nonzero element of the rational function field for
every \(k\ge5\).  Therefore, on every coordinate-minor chart, its expression
can be written as \(P/Q\) with \(Q\ne0\) on the chart and with \(P\) not the
zero polynomial.  The nonzero locus is the complement of the algebraic zero
set \(\{P=0\}\) inside that chart, so it is open dense there.  These chartwise
open dense loci agree on overlaps because they describe the same
basis-independent determinant.  The coordinate-minor charts cover \(D_k\), so
the nonzero determinant locus is open dense in \(D_k\).
\end{proof}
\end{unverified}

\begin{unverified}
\begin{fact}[Generic finite-orbit regularity]%
\label{fact:fg-finite-orbit-regular-open-dense}
Fix a facet count \(F\).  For every simple cyclic word \(w\) of length \(k\) at
least five, pull back the primitive-denominator domain \(D_k\) and the
determinant condition \(\det(I-M_w)\ne0\) to \((\mathbb R^4)^F\) by selecting
the dual vertices listed in \(w\).  On the common intersection of these
pulled-back primitive-denominator domains, the simultaneous determinant
condition for all such words is open dense.

For any flow-graph nondegenerate facet presentation in
Definition~\ref{def:fg-nondegenerate-facet-presentation}, a word that survives
the positive transition-sign test already has all primitive denominators
nonzero.  Therefore, if every surviving length-at-least-five word also
satisfies the corresponding determinant nonvanishing condition, then the
exhaustive run is finite-orbit regular.
\end{fact}

\begin{proof}
There are only finitely many simple cyclic words for a fixed facet count
\(F\).  Fix one word \(w=(s_1,\ldots,s_k)\) with \(k\ge5\).  Let
\[
  \pi_w:(\mathbb R^4)^F\to(\mathbb R^4)^k,\qquad
  (a_1,\ldots,a_F)\mapsto(a_{s_1},\ldots,a_{s_k})
\]
be the coordinate projection for \(w\).  The pulled-back denominator domain is
\(\pi_w^{-1}(D_k)\).

Lemma~\ref{lem:fg-long-word-determinants-open-dense} says that, on \(D_k\),
\(\det(I-M_w)\) is a rational function that is not identically zero.  Its
pullback along \(\pi_w\) is again not identically zero on
\(\pi_w^{-1}(D_k)\): choose a tuple in \(D_k\) where the determinant is
nonzero and extend the listed vectors by arbitrary values for the remaining
facet indices.  Hence the pulled-back nonzero locus is the complement of a
proper algebraic zero set inside \(\pi_w^{-1}(D_k)\), and is open dense there.

Intersect these open dense conditions over the finitely many words of length
at least five.  This gives an open dense set on the common intersection of the
pulled-back primitive-denominator domains.

For a flow-graph nondegenerate presentation, the linear-independence part of
Definition~\ref{def:fg-nondegenerate-facet-presentation} is the hypothesis used
by Lemma~\ref{lem:fg-short-words-no-positive-orbit} for words of length at
most four.  For a length-at-least-five word that survives the transition-sign
test, each primitive denominator is one of the strict inequalities in the
positive transition-sign condition for an adjacent triple, hence is nonzero.
Therefore the determinant condition is defined for every surviving
length-at-least-five word.  If those determinant conditions are nonzero, then
in particular they are nonzero for every surviving word whose closed search
domain is nonempty.  This is exactly finite-orbit regularity.
\end{proof}
\end{unverified}

\begin{unverified}
\begin{proposition}[Presentation-chamber genericity]%
\label{prop:fg-presentation-chamber-genericity}
Fix a facet count \(F\).  Let
\[
  \mathcal C\subset(\mathbb R^4)^F
\]
be a nonempty Euclidean-open set of ordered normalized dual-row lists
\(A=(a_1,\ldots,a_F)\) such that every \(A\in\mathcal C\) defines a bounded
full-dimensional polytope \(K(A)\) with \(0\in\operatorname{int}K(A)\), and
all \(F\) listed inequalities are irredundant.  Then the subset of
\(\mathcal C\) on which
\begin{enumerate}
\item \(\omega_0(a_i,a_j)\ne0\) for every distinct pair \(i\ne j\);
\item every subset of at most four listed facet normals is linearly
  independent; and
\item every simple cyclic word of length at least five which survives the
  flow-graph transition tests has \(\det(I-M_w)\ne0\)
\end{enumerate}
contains a Euclidean-open dense subset of \(\mathcal C\).  On that open dense
subset, Algorithm~\ref{alg:fg-real-exhaustive-search} returns
\(c_{\mathrm{EHZ}}(K(A))\).
\end{proposition}

\begin{proof}
The first two conditions are finite algebraic nonvanishing conditions on
\((\mathbb R^4)^F\): the functions are the bilinear pairings
\(\omega_0(a_i,a_j)\) and the \(r\times r\) minors of each \(r\)-element
sublist for \(1\le r\le4\).  Their nonzero loci are Euclidean open.  They are
dense in \(\mathcal C\) provided the corresponding polynomials do not vanish
identically on \(\mathcal C\); for a Euclidean-open nonempty chamber this
follows because a nonzero polynomial cannot vanish on an open set.

For the determinant conditions, apply
Fact~\ref{fact:fg-finite-orbit-regular-open-dense} to the finite set of simple
cyclic words with \(F\) labels.  The all-pair nonzero-\(\omega_0\) condition
puts every chamber point in the common primitive-denominator domain needed for
these determinant rational functions, whether or not the corresponding facet
pairs are geometrically adjacent.  The proof of
Fact~\ref{fact:fg-finite-orbit-regular-open-dense} shows that each determinant
nonvanishing condition is the complement of a proper algebraic zero set on
that domain.  Restricting finitely many such open dense conditions to the
Euclidean-open chamber \(\mathcal C\) remains open dense.

The resulting open dense subset of \(\mathcal C\) consists of flow-graph
nondegenerate facet presentations in the sense needed by
Definition~\ref{def:fg-nondegenerate-facet-presentation}, because all-pair
nonzero-\(\omega_0\) implies the nonempty-pair condition stated there and the
linear-independence condition is included explicitly.  The determinant
conditions give finite-orbit regularity by
Fact~\ref{fact:fg-finite-orbit-regular-open-dense}.  Hence the capacity
statement is exactly Theorem~\ref{thm:fg-real-capacity-correctness}.
\end{proof}
\end{unverified}

\begin{unverified}
\begin{proposition}[Correctness of idealized exhaustive FG search]%
\label{prop:fg-real-search-correctness}
For a polytope with a flow-graph nondegenerate facet presentation whose
exhaustive flow-graph tube run is finite-orbit regular,
Algorithm~\ref{alg:fg-real-exhaustive-search} either reports that no strict
positive-time simple closed Reeb trajectory represented by a flow-graph tube
word exists, or returns the minimum action among all such represented
trajectories and returns exactly those represented trajectories with action at
most \(c_{\mathrm{FG}}+\delta\), where \(c_{\mathrm{FG}}\) is that minimum
action.
In particular, under the finite-orbit regularity hypothesis the algorithm does
not stop at the long-word singular rejection branch.
\end{proposition}

\begin{proof}
There are finitely many simple cyclic words, so the enumeration terminates.
Every enumerated word is built from primitive triples, and
Lemmas~\ref{lem:fg-primitive-tubes-affine} and~\ref{lem:fg-tube-gluing} show
that the glued closed search data represents exactly the closed
nonnegative-time candidates for that word.  Lemma~\ref{lem:fg-closed-tube-fixed-points}
identifies closed candidates with fixed points of the affine tube map and
identifies strict Reeb trajectories with the fixed points whose reconstructed
segment times are all positive.

If a word contains an adjacent facet pair with empty intersection, then a
strict positive-time representative cannot use that pair because its breakpoint
would lie in that intersection.  If a word contains a triple without positive
flow-graph transition sign, then a strict positive-time representative cannot
use that triple: by
Lemma~\ref{lem:fg-local-transition-regularity-positive-sign}, the two physical
facet transitions adjacent to the middle facet would force the two strict signs
in Definition~\ref{def:fg-transition-sign}.  Hence skipping such a word does
not remove a represented strict trajectory.

For words of length at most four, Lemma~\ref{lem:fg-short-words-no-positive-orbit}
excludes strict positive-time closed trajectories.  For words of length at
least five, finite-orbit regularity makes every relevant fixed-point equation
nonsingular.  Therefore each represented strict trajectory is either recorded
for its cyclic word or correctly absent because its closed search domain is
empty or its fixed point fails the strict segment-time test.  If no strict
trajectory was recorded, the reported absence is correct.  Otherwise taking the
minimum action and applying the final threshold filter gives exactly the stated
output.
\end{proof}
\end{unverified}

\begin{unverified}
\begin{theorem}[Conditional capacity correctness of idealized FG search]%
\label{thm:fg-real-capacity-correctness}
Let \(K\subset\mathbb R^4\) be a full-dimensional convex polytope with
\(0\in\operatorname{int}K\).  Fix a flow-graph nondegenerate facet
presentation
\[
  K=\{x\in\mathbb R^4:\langle a_i,x\rangle\le1,\ i=1,\ldots,F\}.
\]
If the exhaustive flow-graph tube run for this presentation is finite-orbit
regular, then Algorithm~\ref{alg:fg-real-exhaustive-search}, run on this
presentation, returns
\[
  c_{\mathrm{FG}}=c_{\mathrm{EHZ}}(K).
\]
With threshold \(\delta\ge0\), it also returns every represented strict
positive-time simple closed Reeb trajectory with action at most
\(c_{\mathrm{EHZ}}(K)+\delta\).
\end{theorem}

\begin{proof}
By Fact~\ref{fact:fg-simple-capacity-minimizer}, \(c_{\mathrm{EHZ}}(K)\) is
achieved by a minimum-action simple closed Reeb orbit.  Here simple is used in
the active-word sense from the thesis simple-minimizer theorem: the orbit has
a pure facet-direction description in which each used facet direction occurs
in one connected positive-duration interval.  Use only those facets, and write
their cyclic order as the support word of the orbit.  This support word is a
simple cyclic facet word in the sense of
Definition~\ref{def:fg-simple-cyclic-word}.  The positive dwell times and
Lemma~\ref{lem:fg-local-transition-regularity-positive-sign} force the
positive transition signs of every adjacent triple.  Hence
Lemma~\ref{lem:fg-simple-support-orbit-represented} represents this orbit as
a strict fixed point of the closed tube search data for that word, with the
same action.

Proposition~\ref{prop:fg-real-search-correctness} therefore records at least
one trajectory with action \(c_{\mathrm{EHZ}}(K)\).  It records only strict
closed Reeb trajectories, so none of its recorded actions can be less than
\(c_{\mathrm{EHZ}}(K)\).  Hence the reported minimum is exactly
\(c_{\mathrm{EHZ}}(K)\).  The threshold statement follows from the final
filter in Proposition~\ref{prop:fg-real-search-correctness}.
\end{proof}
\end{unverified}

\begin{unverified}
\begin{corollary}[Determinant-generic capacity correctness of idealized FG search]%
\label{cor:fg-real-capacity-correctness-generic}
Fix a facet count \(F\).  Let \(G_F\) be the intersection of the
primitive-denominator domain and determinant-nonzero conditions from
Fact~\ref{fact:fg-finite-orbit-regular-open-dense}.  If a flow-graph
nondegenerate facet presentation with \(F\) listed facets has its listed dual
vertices in \(G_F\), then its exhaustive flow-graph tube run is
finite-orbit regular.  Therefore
Algorithm~\ref{alg:fg-real-exhaustive-search}, run on that presentation,
returns \(c_{\mathrm{EHZ}}(K)\).
The term determinant-generic is relative to this common
primitive-denominator domain.  Proposition~\ref{prop:fg-presentation-chamber-genericity}
is the corresponding chamber-relative statement for fixed ordered normalized
irredundant \(H\)-presentation chambers.
\end{corollary}

\begin{proof}
Fact~\ref{fact:fg-finite-orbit-regular-open-dense} gives finite-orbit
regularity for the exhaustive run.  The capacity statement is then exactly
Theorem~\ref{thm:fg-real-capacity-correctness}.
\end{proof}
\end{unverified}

\begin{remark}[What this formal section does not yet cover]%
\label{rem:fg-real-missing-work}
This section is the idealized real-number algorithm surface.  It does not
prove that the Rust implementation accepts exactly this class, and it does
not analyze f64 rounding.  It also does not cover the exact Rust singular
fixed-set classifier: the theorem above uses finite-orbit regularity to avoid
the relevant singular fixed-point equations, while Rust can classify singular
fixed sets and treat nonpositive-action singular fixed sets as no-orbit
outcomes.  That runtime boundary needs either a separate proof or a
theorem-facing wrapper that rejects singular fixed maps.

The correctness theorem uses the nonempty facet-pair nonzero-\(\omega_0\)
condition from
Definition~\ref{def:fg-nondegenerate-facet-presentation}.  This is the
mathematical boundary currently targeted by the exact Rust rejection logic and
is stronger than the local physical-transition condition needed in the proof of
Lemma~\ref{lem:fg-local-transition-regularity-positive-sign}.

Proposition~\ref{prop:fg-presentation-chamber-genericity} states the
relative-genericity upgrade for a fixed ordered normalized irredundant
\(H\)-presentation chamber that is Euclidean open in the listed dual-row
coordinates.  The chamber model and the choice to impose all-pair
nonzero-\(\omega_0\), rather than only the nonempty-pair condition from
Definition~\ref{def:fg-nondegenerate-facet-presentation}, were reviewed by
Jörn on 2026-07-01.  The all-pair condition is deliberately stronger than
necessary, but acceptable for the thesis genericity statement.

Arbitrary bounded-irredundant input validation, broader Rust correspondence,
and f64 analysis are separate implementation and numerical-analysis
obligations.

The Rust exact search now mirrors the short-word branch of
Algorithm~\ref{alg:fg-real-exhaustive-search}: it validates linear
independence for every subset of at most four listed facet normals, rejects the
input if this theorem-facing contract fails, and skips simple words of length
at most four as no-orbit outcomes.  The closed-word resolver still retains the
structural length-three zero-time diagnostic, because it is useful for tests
and for understanding why short singular fixed sets should not be treated as
capacity candidates.
\end{remark}
